\documentclass[border=10pt]{standalone}
\usepackage[utf8]{inputenc}
\usepackage[spanish]{babel}
\usepackage{tikz}
\usetikzlibrary{shapes.geometric, arrows.meta, positioning, calc}

\begin{document}

\tikzset{
    proceso/.style={
        rectangle,
        rounded corners=3pt,
        draw=cyan!60!black,
        fill=cyan!8,
        text width=3.8cm,
        minimum height=1.4cm,
        align=center,
        font=\small
    },
    resultado/.style={
        rectangle,
        rounded corners=3pt,
        draw=cyan!60!black,
        fill=cyan!8,
        text width=2.2cm,
        minimum height=1.4cm,
        align=center,
        font=\small
    },
    paso/.style={
        circle,
        fill=orange,
        text=white,
        font=\bfseries\small,
        minimum size=0.65cm,
        inner sep=0pt
    },
    pasopeq/.style={
        circle,
        fill=orange,
        text=white,
        font=\bfseries\tiny,
        minimum size=0.4cm,
        inner sep=0pt
    },
    flecha/.style={
        ->,
        >=Stealth,
        thick,
        cyan!60!black
    }
}

\begin{tikzpicture}[node distance=0.6cm and 1.2cm]

% === FILA SUPERIOR (Con reserva) ===
\node[paso] (p1a) {1};
\node[proceso, right=0.3cm of p1a] (caja1a) {Sumar la cantidad de personas que entraron con reserva durante la semana.};

\node[paso, right=1.2cm of caja1a] (p2a) {2};
\node[proceso, right=0.3cm of p2a, text width=4cm] (caja2a) {Multiplicar la cantidad obtenida en el paso \raisebox{-0.1em}{\tikz{\node[pasopeq]{1};}} por 22,5.};

% === FILA INFERIOR (Sin reserva) ===
\node[proceso, below=1.2cm of caja1a] (caja1b) {Sumar la cantidad de personas que entraron sin reserva durante la semana.};
\node[paso, left=0.3cm of caja1b] (p1b) {1};

\node[proceso, right=1.2cm of caja1b, text width=4cm] (caja2b) {Multiplicar la cantidad obtenida en el paso \raisebox{-0.1em}{\tikz{\node[pasopeq]{1};}} por 17.};
\node[paso, left=0.3cm of caja2b] (p2b) {2};

% === PASO 3 (Resultado) ===
% Calcular posición central entre las dos filas
\coordinate (medio) at ($(caja2a.east)!0.5!(caja2b.east)$);
\node[paso, right=2cm of medio] (p3) {3};
\node[resultado, right=0.3cm of p3] (caja3) {Comparar los resultados};

% === FLECHAS ===
% Flechas de paso 1 a paso 2
\draw[flecha] (caja1a.east) -- (p2a.west);
\draw[flecha] (caja1b.east) -- (p2b.west);

% Flechas de paso 2 a paso 3 (convergentes)
\draw[flecha] (caja2a.east) -- ++(0.6,0) |- (caja3.west);
\draw[flecha] (caja2b.east) -- ++(0.6,0) |- (caja3.west);

\end{tikzpicture}

\end{document}
